(1) 热容为常数 $C>0$ 的有限热源初温 $T_i>T_0$，环境恒温 $T_0$。可逆热机从热源取热并向环境放热，直到热源温度降到 $T_0$。求微小过程的功。
(2) 积分求整个过程中最多能获得的功 $W_{\max}$。
(3) 若直接将热源与环境接触而不作功，求不可逆熵产生并与最大功比较。
(4) 对小温差 $T_i=T_0+\delta T$ 展开最大功至二阶，并解释为何不是 $C\delta T$。
